\documentclass[10pt]{beamer}

\usepackage[backend=bibtex,firstinits=true,style=verbose-inote,citestyle=authortitle]{biblatex}
\usepackage{bm}
\usepackage{graphicx}
\usepackage{subcaption}
\usepackage{makecell}
\newcommand{\expect}[2][]{
\ifthenelse{\equal{#1}{}}{
\mathbb{E}\left[#2\right]
}{
\underset{#1}{\mathbb{E}}\left[#2\right]
}}

\newcommand{\cov}[2][]{
\ifthenelse{\equal{#1}{}}{
\text{Cov}\left[#2\right]
}{
\underset{#1}{\text{Cov}}\left[#2\right]
}}


\newcommand{\var}[2][]{
\ifthenelse{\equal{#1}{}}{
\text{Var}[#2]
}{
\underset{#1}{\text{Var}}[#2]
}}

\newcommand{\loss}[2][]{
\ifthenelse{\equal{#1}{}}{
\mathcal{L}(#2)
}{
\mathcal{L}_{#1}(#2)
}}

\newcommand{\kl}[2]{
\text{D}_\text{KL}[#1 \parallel #2]
}

\newcommand{\R}{\mathbb{R}}
%\newcommand{\Prob}{\mathbb{P}}

\newcommand{\1}[1]{\mathds{1}\{#1\}}


%\usecolortheme{dolphin}
\setbeamertemplate{navigation symbols}{}
\setbeamertemplate{section in toc}{\inserttocsectionnumber.~\inserttocsection}
%\addbibresource{style-transfer.bib}

\title{Information Bottleneck}
\subtitle{(Mostly in Deep Learning)}
\author{Ivan Skorokhodov}
%\date{}
%\logo{\includegraphics[height=1cm]{images/ipavlov-logo.png}}

\newcommand{\citepaper}[1]{\citetitle{#1} by \citeauthor{#1}}

%\graphicspath{{./images}}

%\usetheme{lucid}
\begin{document}


\begin{frame}
    \titlepage
\end{frame}


%\begin{frame}
%    \frametitle{Contents}
%    \tableofcontents
%\end{frame}


% http://www.stat.cmu.edu/~larry/=sml/Concentration.pdf
% https://arxiv.org/abs/1904.03743
% https://arxiv.org/abs/1611.01353
% https://reader.elsevier.com/reader/sd/pii/S030439751000201X?token=453A6BFDA0F0963DECE07864F11DCABCF47BD023F089FED6D9168E82AD34920E1D53855647EFB60010239012DD6B28E3
% https://arxiv.org/abs/1612.00410
% https://arxiv.org/abs/1711.00464
% https://arxiv.org/abs/1705.02436
% https://openreview.net/forum?id=SJePKo5HdV
% https://arxiv.org/abs/1808.07593
% https://arxiv.org/abs/1901.10799
% https://papers.nips.cc/paper/8120-theoretical-linear-convergence-of-unfolded-ista-and-its-practical-weights-and-thresholds
% https://openreview.net/forum?id=Bklr3j0cKX
% https://openreview.net/forum?id=BkedwoC5t7
% https://openreview.net/forum?id=HyxPx3R9tm
% https://www.sciencedirect.com/science/article/pii/S0893608019301352
% https://arxiv.org/abs/1906.04109
% https://arxiv.org/abs/1710.05233
% https://arxiv.org/abs/1802.09766
% https://arxiv.org/abs/1808.07593
% https://arxiv.org/abs/1801.04062
% https://physics.princeton.edu/archives/theses/lib/upload/Strouse-thesis.pdf
% http://pillowlab.princeton.edu/teaching/statneuro2018/slides/notes08_infotheory.pdf

\section{Style Transfer not in Text}
\subsection{Image Style Transfer}


\begin{frame}
\frametitle{Plan}
%\framesubtitle{Overview}

\begin{enumerate}
    \item Origins and intuitive understanding
    \item IB is not a hidden variable model (no generative assumption)
    \item It usually assumes that joint is known
    \item Connection to minimal sufficient statistics
    \item Connection to VAE/variational inference in general
    \item Trivial solution for $\beta \leq 1$.
    \item Problems in deterministic scenario
    \item Implicit optimization in DL: SZT experiments and Saxe's critics
    \item Explicit optimization in DL: TODO
    \item Why IB? --- Generalization bounds
    \item Why IB? --- Bias-variance tradeoff
    \item MI estimators
\end{enumerate}
\end{frame}

\begin{frame}
\frametitle{Origins and intuitive understanding}
\begin{itemize}
    \item TODO: do we always can compute $z$ from a new $x$, i.e. $W$ is always invertible?
    \item PCA solves:
\begin{equation*}
\min_{W,Z} \| WZ - X\|_F
\end{equation*}
i.e. we try to find linear mapping $W$ and latent codes $Z$ such that $X$ is ``reconstructed well'' from $Z$.
\item IB solves
\begin{equation*}
\min_{Z} I(X:Z) - \beta I(Z:Y)
\end{equation*}
i.e. we try to find such $Z$ that $Y$ is ``reconstructed well'' and \textit{we do not care about $X$}.
\item Q1: Why should we care about minimizing $I(X:Z)$ if we only care about reconstructing $Y$?
\item A1: $I(X:Z) = H(Z) - H(Z|X)$, so $H(Z)$ is minimized (a good property)
\item Q2: But $H(Z|X)$ is maximized: why?
\item A2: The reason is subtle, we can do well without that
\end{itemize}
\end{frame}

\begin{frame}
\begin{equation*}
\begin{split}
\kl{p(w)}{\tilde{q}(w)} &= \int \log\frac{p(w)}{\tilde{q}(w)} p(w) dw \\
&= \int \log\frac{p(w)}{\tilde{p}(w)} p(w) dw + \int \log\frac{\tilde p(w)}{\tilde q(w)} p(w) dw
\\
&= \kl{p(w)}{\tilde p(w)} + \int \sum_{i=1}^d \log \frac{\tilde p(w_i)}{\tilde q(w_i)} p(w) dw
\\
&= \kl{p(w)}{\tilde p(w)} + \sum_{i=1}^d \kl{\tilde p(w_i)}{\tilde q(w_i)}
\\
&= \kl{p(w)}{\tilde p(w)} + \kl{\tilde p(w)}{\tilde q(w)}
\end{split}
\end{equation*}
\end{frame}

\end{document}
